\section{Mathematical Model}

\begin{longtable}{|p{0.20\textwidth}|p{0.79\textwidth}|}
\caption{Model nomenclature}
\label{tab:nomenclature}\\
\hline
\multicolumn{2}{|l|}{\textbf{Indices and Sets}}\\ \hline
$V$          & Set of nodes representing manufacturing plants, $i,j \in V$ \\
$E$          & Set of directed edges representing potential transport lanes,
               $E \subseteq V \times V$ \\
$C$          & Set of RTI types, $c \in C$ \textit{(to be incorporated)} \\
\hline
\multicolumn{2}{|l|}{\textbf{Parameters}}\\ \hline
$\mu_{ij}$        & Daily average number of full RTIs shipped on lane $(i,j)$ \\
$ric_{ij}$        & Fraction of full shipments on lane $(i,j)$ that require a
                    return inlay \\
$c_{ij}$          & Transport cost per unit volume on lane $(i,j)$ \\
$f_{ij}$          & Fixed cost per shipment on lane $(i,j)$ \\
$v^{f}$, $v^{e}$  & Volume occupied by one full and one empty RTI,
                    respectively \\
$l_{ij}$          & Transit time on lane $(i,j)$, in days \\
$\alpha^{f}$, $\alpha^{e}$ & Maintenance loss factor for full and empty RTIs,
                    respectively \\
$ssf^{i}_{ij}$, $ssf^{j}_{ij}$ & Safety stock of full RTIs held at the origin
                    and destination node of lane $(i,j)$ \\
$sse^{i}_{ij}$, $sse^{j}_{ij}$ & Safety stock of empty RTIs held at the origin
                    and destination node of lane $(i,j)$ \\
$s^{\text{RTI}}$  & Initial RTI pool available before optimisation \\
$p^{\text{RTI}}$  & Unit purchase cost of one RTI \\
$\omega_{ij}$     & Minimum shipment volume on lane $(i,j)$ \\
$\delta_{ij}$     & Maximum shipment volume on lane $(i,j)$ \\
$M$               & Sufficiently large constant (big-$M$) \\
\hline
\multicolumn{2}{|l|}{\textbf{Decision Variables}}\\ \hline
\multicolumn{2}{|l|}{\textit{Binary variables}}\\ \hline
$X_{ij}$     & 1 if lane $(i,j)$ is active for any RTI flow, 0 otherwise \\
$XF_{ij}$    & 1 if lane $(i,j)$ carries full RTIs, 0 otherwise \\
$XI_{ij}$    & 1 if lane $(i,j)$ carries return inlays, 0 otherwise \\
$XE_{ij}$    & 1 if lane $(i,j)$ carries empty RTIs, 0 otherwise \\
\hline
\multicolumn{2}{|l|}{\textit{Daily flow variables}}\\ \hline
$NF_{ij}$    & Daily average number of full RTIs sent from $i$ to $j$ \\
$NI_{ij}$    & Daily average number of return inlays sent from $i$ to $j$ \\
$NE_{ij}$    & Daily average number of empty RTIs sent from $i$ to $j$ \\
\hline
\multicolumn{2}{|l|}{\textit{Daily volume variables}}\\ \hline
$VF_{ij}$    & Daily volume of full RTIs on lane $(i,j)$ \\
$VI_{ij}$    & Daily volume of return inlays on lane $(i,j)$ \\
$VE_{ij}$    & Daily volume of empty RTIs on lane $(i,j)$ \\
\hline
\multicolumn{2}{|l|}{\textit{Per-shipment variables}}\\ \hline
$q_{ij}$     & Queue time on lane $(i,j)$: days between consecutive
               dispatches (continuous, $q_{ij} \geq 1$) \\
$VS_{ij}$    & Total volume per shipment on lane $(i,j)$
               [$= (VF_{ij}+VI_{ij}+VE_{ij})\cdot q_{ij}$,\ \textit{bilinear}] \\
$NSF_{ij}$, $NSI_{ij}$, $NSE_{ij}$ & Number of full, inlay and empty RTIs
               per shipment on lane $(i,j)$
               [$= NF_{ij}\cdot q_{ij}$, etc.,\ \textit{bilinear}] \\
\hline
\multicolumn{2}{|l|}{\textit{Total RTI pool variables}}\\ \hline
$TNCF_{ij}$  & Total full RTIs allocated to lane $(i,j)$, accounting for
               in-transit stock, safety stock, and maintenance losses \\
$TNCI_{ij}$  & Total inlay RTIs allocated to lane $(i,j)$ \\
$TNCE_{ij}$  & Total empty RTIs allocated to lane $(i,j)$, accounting for
               in-transit stock, safety stock, and maintenance losses \\
$P$          & Number of additional RTIs to purchase ($P \geq 0$) \\
\hline
\multicolumn{2}{|l|}{\textit{Cost variables}}\\ \hline
$CF_{ij}$    & Daily transport cost of full RTIs on lane $(i,j)$ \\
$CI_{ij}$    & Daily transport cost of return inlays on lane $(i,j)$ \\
$CE_{ij}$    & Daily transport cost of empty RTIs on lane $(i,j)$ \\
$TC_{ij}$    & Total daily transport cost on lane $(i,j)$ \\
$TCS_{ij}$   & Transport cost per shipment on lane $(i,j)$
               [$= VS_{ij}\cdot c_{ij} + X_{ij}\cdot f_{ij}$] \\
\hline
\end{longtable}

% ─────────────────────────────────────────────────────────────────
\paragraph{Objective function.}
Minimise the total daily transport cost across all lanes plus the amortised
purchase cost of any additional RTIs required:
%
\begin{align}
  \min \quad \sum_{(i,j)\in E} TC_{ij} \;+\; p^{\text{RTI}} \cdot P
  \label{eq:objective}
\end{align}

% ─────────────────────────────────────────────────────────────────
\paragraph{Flow conservation.}
The total number of RTIs entering each node equals the total number leaving it,
ensuring a closed-loop network:
%
\begin{align}
  \sum_{j \in V}\!\bigl(NE_{ij}+NF_{ij}+NI_{ij}\bigr)
  &= \sum_{j \in V}\!\bigl(NE_{ji}+NF_{ji}+NI_{ji}\bigr)
  && \forall\, i \in V
  \label{eq:flow-conservation}
\end{align}

% ─────────────────────────────────────────────────────────────────
\paragraph{Fixed full-RTI and inlay flows.}
Full RTI flows are fixed by production plans. Return inlays are sent back on
the reverse lane in proportion to the full shipments that required them:
%
\begin{align}
  NF_{ij} &= \mu_{ij}
    && \forall\,(i,j)\in E
    \label{eq:full-fixed}\\
  NI_{ji} &= ric_{ij}\cdot\mu_{ij}
    && \forall\,(i,j)\in E
    \label{eq:inlay-fixed}
\end{align}

% ─────────────────────────────────────────────────────────────────
\paragraph{Lane activation.}
A lane is active if and only if it carries at least one RTI flow type:
%
\begin{align}
  X_{ij} &\geq XF_{ij}, \quad X_{ij} \geq XI_{ij}, \quad X_{ij} \geq XE_{ij}
    && \forall\,(i,j)\in E
    \label{eq:activation-lb}\\
  X_{ij} &\leq XF_{ij} + XI_{ij} + XE_{ij}
    && \forall\,(i,j)\in E
    \label{eq:activation-ub}
\end{align}

% ─────────────────────────────────────────────────────────────────
\paragraph{Flow--binary consistency.}
Each binary is active if and only if a positive flow of the corresponding type
is present:
%
\begin{align}
  NF_{ij} &\leq M\cdot XF_{ij}, \quad XF_{ij} \leq NF_{ij}
    && \forall\,(i,j)\in E
    \label{eq:binary-full}\\
  NI_{ij} &\leq M\cdot XI_{ij}, \quad XI_{ij} \leq NI_{ij}
    && \forall\,(i,j)\in E
    \label{eq:binary-inlay}\\
  NE_{ij} &\leq M\cdot XE_{ij}, \quad XE_{ij} \leq NE_{ij}
    && \forall\,(i,j)\in E
    \label{eq:binary-empty}
\end{align}

% ─────────────────────────────────────────────────────────────────
\paragraph{Daily volumes.}
The daily volume on each lane is determined by the number of RTIs and their
individual volumes:
%
\begin{align}
  VF_{ij} &= NF_{ij}\cdot v^{f}
    && \forall\,(i,j)\in E
    \label{eq:vol-full}\\
  VI_{ij} &= NI_{ij}\cdot v^{f}
    && \forall\,(i,j)\in E
    \label{eq:vol-inlay}\\
  VE_{ij} &= NE_{ij}\cdot v^{e}
    && \forall\,(i,j)\in E
    \label{eq:vol-empty}
\end{align}

% ─────────────────────────────────────────────────────────────────
\paragraph{Per-shipment quantities.}
The number of RTIs and total volume accumulated between consecutive dispatches
are proportional to the queue time. These constraints are \emph{bilinear} and
will be linearised in Section~\ref{sec:linearisation}:
%
\begin{align}
  NSF_{ij} &= NF_{ij}\cdot q_{ij}
    && \forall\,(i,j)\in E
    \label{eq:ship-full}\\
  NSI_{ij} &= NI_{ij}\cdot q_{ij}
    && \forall\,(i,j)\in E
    \label{eq:ship-inlay}\\
  NSE_{ij} &= NE_{ij}\cdot q_{ij}
    && \forall\,(i,j)\in E
    \label{eq:ship-empty}\\
  VS_{ij}  &= \bigl(VF_{ij}+VI_{ij}+VE_{ij}\bigr)\cdot q_{ij}
    && \forall\,(i,j)\in E
    \label{eq:ship-vol}
\end{align}

% ─────────────────────────────────────────────────────────────────
\paragraph{Shipment volume bounds.}
Each shipment must respect a minimum load and a maximum vehicle capacity:
%
\begin{align}
  \omega_{ij}\cdot X_{ij} &\leq VS_{ij} \leq \delta_{ij}\cdot X_{ij}
    && \forall\,(i,j)\in E
    \label{eq:vol-bounds}
\end{align}

% ─────────────────────────────────────────────────────────────────
\paragraph{Total RTI pool requirements.}
The total number of RTIs allocated to each lane accounts for in-transit stock,
safety stock at both endpoints, and a maintenance loss factor:
%
\begin{align}
  TNCF_{ij} &= (1+\alpha^{f})\bigl(NSF_{ij} + l_{ij}\cdot NF_{ij}
               + XF_{ij}(ssf^{i}_{ij}+ssf^{j}_{ij})\bigr)
    && \forall\,(i,j)\in E
    \label{eq:pool-full}\\
  TNCE_{ij} &= (1+\alpha^{e})\bigl(NSE_{ij} + l_{ij}\cdot NE_{ij}
               + XE_{ij}(sse^{i}_{ij}+sse^{j}_{ij})\bigr)
    && \forall\,(i,j)\in E
    \label{eq:pool-empty}\\
  TNCI_{ij} &= (1+\alpha^{f})\bigl(NSI_{ij} + l_{ij}\cdot NI_{ij}\bigr)
    && \forall\,(i,j)\in E
    \label{eq:pool-inlay}
\end{align}

\paragraph{RTI purchase.}
Additional RTIs must be purchased to cover any shortfall between the network's
total requirement and the existing pool. If the existing pool is sufficient,
no purchase is needed and the surplus remains idle:
%
\begin{align}
  \sum_{(i,j)\in E}\bigl(TNCF_{ij}+TNCE_{ij}+TNCI_{ij}\bigr)
    &\leq s^{\text{RTI}} + P
    && P \geq 0
  \label{eq:purchase}
\end{align}

% ─────────────────────────────────────────────────────────────────
\paragraph{Daily transport costs.}
The per-shipment cost is the sum of a variable volumetric component and a fixed
dispatch cost. Dividing by the queue time converts it to a daily cost. This
constraint is \emph{bilinear} and will be linearised in
Section~\ref{sec:linearisation}:
%
\begin{align}
  TCS_{ij} &= VS_{ij}\cdot c_{ij} + X_{ij}\cdot f_{ij}
    && \forall\,(i,j)\in E
    \label{eq:cost-per-ship}\\
  TC_{ij}\cdot q_{ij} &= TCS_{ij}
    && \forall\,(i,j)\in E
    \label{eq:cost-daily}
\end{align}
